\documentclass[preview]{standalone}

\usepackage{amsmath}
\usepackage{amssymb}
\usepackage{stellar}
\usepackage{definitions}

\begin{document}

\id{quotient-spaces-dimension}
\genpage

\section{Dimension of quotient spaces}

\begin{snippetproposition}{dimension-of-quotient-space}{Dimension of a quotient space}
    Let \(V\) be a finite-dimensional \vectorspace over a \field
    \(\mathbb{F}\), and let \(W\subseteq V\) be a linear subspace. Then
    \[
        \lineardim(V/W)=\lineardim V-\lineardim W.
    \]
\end{snippetproposition}

\begin{snippetproof}{dimension-of-quotient-space-proof}{dimension-of-quotient-space}{Dimension of a quotient space}
    Write \(\lineardim V=n\) and \(\lineardim W=r\). Let
    \[
        \mathcal{B}_W=\{w_1,\ldots,w_r\}
    \]
    be a \basis of \(W\), and extend it to a \basis
    \[
        \mathcal{B}=\{w_1,\ldots,w_r,v_1,\ldots,v_{n-r}\}
    \]
    of \(V\). We show that
    \[
        \mathcal{B}_{V/W}=\{v_1+W,\ldots,v_{n-r}+W\}
    \]
    is a \basis of \(V/W\).

    Suppose that
    \[
        a_1(v_1+W)+\cdots+a_{n-r}(v_{n-r}+W)=0_{V/W}.
    \]
    Then \(a_1v_1+\cdots+a_{n-r}v_{n-r}\in W\), so there are
    \(b_1,\ldots,b_r\in\mathbb{F}\) such that
    \[
        a_1v_1+\cdots+a_{n-r}v_{n-r}
        =b_1w_1+\cdots+b_rw_r.
    \]
    The \linearlyindependent[linear independence] of \(\mathcal{B}\) gives
    \(a_1=\cdots=a_{n-r}=0\). Thus \(\mathcal{B}_{V/W}\) is
    \linearlyindependent.

    Now let \(v+W\in V/W\). There are scalars
    \(b_1,\ldots,b_r,a_1,\ldots,a_{n-r}\in\mathbb{F}\) such that
    \[
        v=b_1w_1+\cdots+b_rw_r+a_1v_1+\cdots+a_{n-r}v_{n-r}.
    \]
    Since \(b_1w_1+\cdots+b_rw_r\in W\), passage to the quotient yields
    \[
        v+W=a_1(v_1+W)+\cdots+a_{n-r}(v_{n-r}+W).
    \]
    Hence \(\mathcal{B}_{V/W}\) generates \(V/W\). It is therefore a
    \basis with \(n-r\) elements, and
    \(\lineardim(V/W)=n-r\).
\end{snippetproof}

\end{document}
